\documentclass[11pt]{article}

\usepackage[a4paper,margin=1in]{geometry}
\usepackage[T1]{fontenc}
\usepackage[utf8]{inputenc}
\usepackage{lmodern}
\usepackage{microtype}
\usepackage{amsmath,amssymb}
\usepackage{booktabs}
\usepackage{graphicx}
\usepackage{siunitx}
\usepackage{caption}
\usepackage{subcaption}
\usepackage{enumitem}
\usepackage{xcolor}
\usepackage{hyperref}
\usepackage[nameinlink,noabbrev]{cleveref}

\graphicspath{{figures/}}
\sisetup{separate-uncertainty=true}

\hypersetup{
  colorlinks=true,
  linkcolor=blue!55!black,
  citecolor=blue!55!black,
  urlcolor=blue!55!black,
  pdftitle={SiPM Operating-Point Study for the gLOWCOST Detector},
  pdfauthor={Tharindu Hettiarachchi}
}

\newcommand{\vbr}{V_{\mathrm{br}}}
\newcommand{\vbias}{V_{\mathrm{bias}}}
\newcommand{\vov}{V_{\mathrm{OV}}}
\newcommand{\pe}{\mathrm{p.e.}}
\newcommand{\studysection}[1]{\section{#1}\leavevmode\par}
\newcommand{\studysubsection}[1]{\subsection{#1}\leavevmode\par}
\newcommand{\studysubsubsection}[1]{\subsubsection{#1}\leavevmode\par}
\newcommand{\studyfigure}[3]{%
  \begin{figure}[htbp]
    \centering
    \IfFileExists{figures/#1}{%
      \includegraphics[width=0.90\linewidth]{#1}%
    }{%
      \fbox{\parbox[c][50mm][c]{0.86\linewidth}{\centering
      Figure placeholder\\[4pt]\texttt{\detokenize{#1}}}}%
    }
    \caption{#2}
    \label{#3}
  \end{figure}%
}

\title{Characterization and Operating-Point Calibration of SiPMs\\
for the gLOWCOST Detector}
\author{Tharindu Hettiarachchi\\
Department of Physics and Astronomy, Georgia State University}
\date{\today}

\begin{document}

\maketitle

\begin{abstract}
% Add a short statement of the problem, method, principal results, and conclusion.
\end{abstract}

\tableofcontents
\newpage

\studysection{Introduction}
\studysubsection{Motivation}
\studysubsection{The gLOWCOST detector}
\studysubsection{Why the SiPM operating point must be calibrated}
\studysubsection{Objectives of this study}
\studysubsection{Scope of the work}

\studysection{Detector and Instrumentation}
\studysubsection{Hamamatsu S13360-2050VE SiPM}
\studysubsection{GSU gLOWCOST scintillator tiles}
\studysubsection{EPIC reference tiles}
\studysubsection{Wavelength-shifting fibre and optical coupling}
\studysubsection{Analogue readout electronics}
\studysubsubsection{SiPM bias and signal connection}
\studysubsubsection{Amplification stage}
\studysubsubsection{Comparator and FPGA counting path}
\studysubsection{Bias-voltage generation}
\studysubsubsection{MAX1932 high-side supply}
\studysubsubsection{Per-channel DAC adjustment}
\studysubsubsection{Effective SiPM bias}
\studysubsection{PicoScope waveform acquisition}
\studysubsection{Temperature and environmental sensors}
\studysubsection{Channel mapping and detector geometry}

\begin{table}[htbp]
  \centering
  \caption{Channel mapping used for the EPIC-triggered GSU gLOWCOST measurement.}
  \label{tab:channel-map}
  \begin{tabular}{llll}
    \toprule
    Detector element & SiPM label & PCB channel & Scope channel \\
    \midrule
    Top EPIC tile &  & CH0 & C \\
    Top GSU gLOWCOST tile & SiPM 1 (Triangle) & CH2 & A \\
    Bottom GSU gLOWCOST tile & SiPM 2 (Star) & CH3 & B \\
    Bottom EPIC tile &  & CH1 & D \\
    \bottomrule
  \end{tabular}
\end{table}

\studysection{Physical Basis}
\studysubsection{Geiger-mode avalanche operation}
\studysubsection{Breakdown voltage and overvoltage}

\begin{equation}
  \vov = \vbias - \vbr .
  \label{eq:overvoltage}
\end{equation}

\studysubsection{SiPM gain and photoelectron peak spacing}
\studysubsection{Photon-detection efficiency}
\studysubsection{Dark counts}
\studysubsection{Optical crosstalk and afterpulsing}
\studysubsection{Temperature dependence of breakdown voltage}
\studysubsection{Scintillation-light production and collection}
\studysubsection{Zero-event probability}

\begin{equation}
  P(0) = e^{-\mu},
  \qquad
  \mu = -\ln P(0).
  \label{eq:poisson-zero}
\end{equation}

\studysubsection{Dark-window correction}

\begin{equation}
  P_{0,\mathrm{signal}} = \frac{f_{0,\mathrm{signal\ gate}}}
  {f_{0,\mathrm{dark\ gate}}},
  \qquad
  \epsilon_{\mathrm{chain}} = 1-P_{0,\mathrm{signal}}.
  \label{eq:dark-corrected-zero}
\end{equation}

\studysection{Experimental Method}
\studysubsection{SiPM identification and labeling}
\studysubsection{Electrical and optical connections}
\studysubsection{Probe attenuation and oscilloscope configuration}
\studysubsection{Bias-voltage settings}
\studysubsection{Temperature stabilization and recording}
\studysubsection{Light-blocked dark-pulse acquisition}
\studysubsection{Breakdown-voltage scan}
\studysubsubsection{Voltage range and step size}
\studysubsubsection{Events recorded at each voltage}
\studysubsubsection{Bias settling time}
\studysubsubsection{Trigger settings}
\studysubsection{Channel-swap test}
\studysubsection{Independent repeat scan}
\studysubsection{EPIC-tile reference coincidence}
\studysubsection{GSU gLOWCOST zero-event acquisition}
\studysubsection{Control measurements}
\studysubsubsection{High-voltage-off control}
\studysubsubsection{Dark-time control window}
\studysubsubsection{Saturation and overflow checks}

\studysection{Waveform Processing}
\studysubsection{Raw PicoScope data format}
\studysubsection{Probe-attenuation correction}
\studysubsection{Event-by-event baseline estimation}
\studysubsection{Baseline-noise measurement}
\studysubsection{Pulse search window}
\studysubsection{Pulse-height measurement}
\studysubsection{Peak-aligned pulse-area integration}
\studysubsection{Waveform acceptance criteria}
\studysubsection{Clipping and saturation detection}
\studysubsection{Dark and signal integration gates}

\studyfigure{example-waveform.pdf}%
{Example waveform showing the baseline region, pulse-search region, and integration gate.}%
{fig:waveform-example}

\studysection{Photoelectron Spectrum Analysis}
\studysubsection{Construction of pulse-height histograms}
\studysubsection{Construction of pulse-area histograms}
\studysubsection{Statistical bin-width selection}
\studysubsection{Histogram smoothing for candidate discovery}
\studysubsection{Candidate peak detection}
\studysubsection{Local peak fitting}
\studysubsection{Selection of evenly spaced photoelectron populations}
\studysubsection{Pedestal-constrained multi-population model}
\studysubsection{Photoelectron-gap uncertainty}
\studysubsection{Tests for unresolved or misleading spectra}

\studyfigure{example-pe-spectrum.pdf}%
{Example photoelectron spectrum with the fitted population centers.}%
{fig:pe-spectrum}

\studysection{Breakdown-Voltage Determination}
\studysubsection{Photoelectron spacing as a gain proxy}
\studysubsection{Linear spacing model}

\begin{equation}
  \Delta Q = m\vbias+b = m(\vbias-\vbr),
  \qquad
  \vbr=-\frac{b}{m}.
  \label{eq:vbr-fit}
\end{equation}

\studysubsection{Weighted linear regression}
\studysubsection{Bias-point acceptance}
\studysubsection{Temperature normalization}
\studysubsection{Fit uncertainty}
\studysubsection{Voltage-calibration systematic uncertainty}
\studysubsection{Leave-one-out and repeatability checks}

\studyfigure{example-vbr-fit.pdf}%
{Photoelectron spacing as a function of effective bias. The x-intercept gives the breakdown voltage.}%
{fig:vbr-fit}

\studysection{Zero-Event Analysis}
\studysubsection{Selection of EPIC-tile coincidence events}
\studysubsection{Timing-coincidence requirement}
\studysubsection{GSU gLOWCOST signal window}
\studysubsection{Dark-time window}
\studysubsection{Definition of a zero event}
\studysubsection{Correction for dark counts and electronic noise}
\studysubsection{Full-chain detection efficiency}
\studysubsection{Zero-equivalent mean occupancy}
\studysubsection{Statistical interval calculation}
\studysubsection{Accidental-coincidence estimate}
\studysubsection{Interpretation limits for cosmic-ray data}

\studysection{Results}
\studysubsection{Bias-voltage calibration}
\studysubsection{Recorded waveform quality}
\studysubsection{Photoelectron spectra}
\studysubsection{Photoelectron spacing versus bias}
\studysubsection{Breakdown voltages}

\begin{table}[htbp]
  \centering
  \caption{Measured SiPM breakdown voltages and provisional operating points.}
  \label{tab:vbr-results}
  \begin{tabular}{lccc}
    \toprule
    Device & $\vbr$ (V) & Reference temperature ($^\circ$C) & Operating bias (V) \\
    \midrule
    SiPM 1 (Triangle) &  &  &  \\
    SiPM 2 (Star) &  &  &  \\
    \bottomrule
  \end{tabular}
\end{table}

\studysubsection{Repeatability of the breakdown-voltage measurement}
\studysubsection{Channel-swap comparison}
\studysubsection{Temperature dependence}
\studysubsection{Dark-count response at the operating point}
\studysubsection{EPIC-triggered GSU gLOWCOST response}
\studysubsection{Zero-event fractions}
\studysubsection{Full-chain detection efficiencies}
\studysubsection{Comparison of SiPM 1 and SiPM 2}

\studysection{Discussion}
\studysubsection{Physical interpretation of the photoelectron spectra}
\studysubsection{Validity of the linear gain extrapolation}
\studysubsection{Comparison with Hamamatsu specifications}
\studysubsection{Effect of the analogue readout chain}
\studysubsection{Effect of trigger selection}
\studysubsection{Device differences and channel-dependent effects}
\studysubsection{Meaning of equal overvoltage}
\studysubsection{Meaning of matched full-chain efficiency}
\studysubsection{Statistical uncertainty}
\studysubsection{Systematic uncertainty}
\studysubsection{Limitations of the scintillator zero-event measurement}
\studysubsection{Implications for the gLOWCOST detector operating point}

\studysection{Recommended Operating Point}
\studysubsection{Selection criteria}
\studysubsection{Bias setting for each SiPM}
\studysubsection{Temperature-compensation coefficient}
\studysubsection{Detector threshold}
\studysubsection{Expected operating margins}

\studysection{Future Measurements}
\studysubsection{Direct calibration of the delivered bias voltage}
\studysubsection{Controlled-temperature breakdown-voltage scan}
\studysubsection{Pulsed-LED zero-event measurement}
\studysubsection{Relative photon-detection-efficiency matching}
\studysubsection{Dark-count, crosstalk, and afterpulse measurements}
\studysubsection{Long-term detector stability}

\studysection{Conclusion}

\section*{Acknowledgements}
\leavevmode\par
\addcontentsline{toc}{section}{Acknowledgements}

\nocite{hamamatsu_s13360_series}
\bibliographystyle{unsrt}
\bibliography{references}

\appendix

\studysection{Complete Acquisition Settings}
\studysubsection{Breakdown-voltage scans}
\studysubsection{Zero-event acquisition}

\studysection{Bias-Control Codes and Voltage Conversion}

\studysection{Waveform-Analysis Parameters}

\studysection{Peak-Finding Parameters}

\studysection{Additional Spectra and Fit Diagnostics}

\studysection{Rejected and Diagnostic Runs}

\studysection{Data and Code Availability}

\end{document}
